\section{Untruncated optical transport and fractional potentials}
\label{sec:v83-potential}
Fix one of the original exact events. The field, optical density and
coarea kernels in Input~\ref{input:v83-geometric} are used throughout.
No new field is chosen from the unknown roof density. Put $z=|g|$.
For $0<\eta<1/2$ set
\begin{equation}\label{eq:v83-potential-parameters}
 \kappa_\eta=\frac{1-\eta}{2},\qquad Q_\eta=\kappa_\eta^{-1},\qquad
 \mathsf P_{\eta,r}v(t)=\int_0^r s^{-\kappa_\eta}v(t+s)\,ds.
\end{equation}
In particular $2<Q_\eta<4$. The source $v^\eta$ below is the
\emph{actual}, untruncated source with insertion $I_m z^{-\eta}$
from Input~\ref{input:v83-local}; it is not a limit taken after the
collision limit.

Write $E_h^{\varepsilon,\mathrm{gr}}$ for the starting mass of genuine
grazing first exits and $O_h^\varepsilon$ for the positive inverse-twist
loss on survivors. These are disjoint classes of initial states. If
$b>0$, a subscript $<b$ restricts the initial source to $0<z<b$.
A bar denotes the average $r^{-1}\int_0^r\cdot\,dh$.

\begin{theorem}[Optical transport without an initial lower cutoff]
\label{thm:v83-potential-transport}
For almost every $(t,h)$,
\begin{equation}\label{eq:v83-untruncated-transport}
 O_h^\varepsilon(t)+E_h^{\varepsilon,\mathrm{gr}}(t)
                   \le C_\eta\mathsf P_{\eta,h}v^\eta(t).
\end{equation}
The right side is locally integrable in $t$ at every fixed count.
If $0<h\le r$ and $b_r=(b^2+2\Lambda_\circ r)^{1/2}<b_1$, then
\begin{equation}\label{eq:v83-critical-potential}
 (O_h^\varepsilon+E_h^{\varepsilon,\mathrm{gr}})_{<b}(t)
              \le C_\eta\mathsf P_{\eta,r}v^{\eta,b_r}(t).
\end{equation}
The same bounds hold for the $h$-averaged sources. All internal
incidences, all actual words and both section tests are retained.
\end{theorem}
\begin{proof}
Along a regular survivor, the optical sensitivity and radial speed give
\[
 |X\omega|\le C\omega I_m/z,\qquad
 z(\Phi_sx)^2\ge z(x)^2+2\lambda_\circ s.
\]
Consequently
\begin{equation}\label{eq:v83-fractional-drift}
 |X\omega|(\Phi_sx)
 \le C_\eta s^{-\kappa_\eta}
       \omega(\Phi_sx)I_m(\Phi_sx)z(\Phi_sx)^{-\eta}.
\end{equation}
For a surviving state, $[\omega(x)-\omega(\Phi_hx)]_+$ is bounded
by the negative variation up to $h$. At a grazing first exit the
limiting optical density is zero; hence its entire positive starting
optical weight is bounded by negative variation up to the lifetime.
This uses the one-word optical identity only to identify the zero
limit, not to sum wordwise constants.

In an angular flow box the initial fiber density is
$\omega(x)/(4\pi c_*\chi(x))\,d\vartheta$. Multiplying by
$W^\varepsilon(x)\le1$ and using
\[
 \frac{\omega(\Phi_sx)}{\omega(x)}
       =J_s(x)\frac{\chi(\Phi_sx)}{\chi(x)}\le C J_s(x)
\]
transfers the integrand in \eqref{eq:v83-fractional-drift} to a subset
of the complete target fiber. The common-kernel covariance in
Appendix~\ref{app:v83-kernels} justifies this simultaneously for the
disjoint survivor and first-exit selections. Tonelli then proves
\eqref{eq:v83-untruncated-transport}. Nonnegative truncation at fixed
collision count removes all compact-chart restrictions.

For $z(x)<b$ the radial upper bound gives $z(\Phi_sx)<b_r$ before
$h\le r$, proving \eqref{eq:v83-critical-potential}. The kernel and
sources are nonnegative, so averaging in $h$ preserves the bounds.
Finally $\kappa_\eta<1$ and $v^\eta$ is locally integrable with compact
support. Fubini proves local integrability of each displayed potential.
Neither an untruncated first inverse-gradient moment nor a pointwise
bound for $v^\eta$ has been used.
\end{proof}

For $Q>1$, a measurable interval $J$ and $f\ge0$, define
\begin{equation}\label{eq:v83-lorentz-norm}
 \mathcal N_Q(f;J)=\sup_{E\subset J,\ 0<|E|<\infty}
          |E|^{-1+1/Q}\int_E f(t)\,dt,
\end{equation}
where the supremum is over Lebesgue-measurable sets. This is a
specific weak-$L^Q$ norm; the definition fixes its normalization.

\begin{lemma}[A localized potential bound]
\label{lem:v83-potential-norm}
For every nonnegative locally integrable $v$, every $r>0$, and any
interval $J$, write $J_r=J+(0,r)$. Then
\begin{equation}\label{eq:v83-potential-norm}
 \mathcal N_{Q_\eta}(\mathsf P_{\eta,r}v;J)
               \le\frac1{1-\kappa_\eta}\int_{J_r}v(y)\,dy.
\end{equation}
The inequality also holds for extended nonnegative quantities.
\end{lemma}
\begin{proof}
For fixed $y$ and a measurable set $E$ of length $a$, the decreasing
kernel has
\[
 \int_E (y-t)^{-\kappa_\eta}\one_{\{0<y-t<r\}}\,dt
 \le \int_0^{\min(a,r)}s^{-\kappa_\eta}\,ds
 \le\frac{a^{1-\kappa_\eta}}{1-\kappa_\eta}.
\]
The first inequality follows directly by integrating the measures of
the kernel's superlevel intervals. Integrate against $v(y)$ and use
Tonelli. Only $y\in J_r$ contributes. Divide by
$a^{1-\kappa_\eta}$ and take the supremum over $E$.
\end{proof}

Set $D_r^\varepsilon=\bar O_r^\varepsilon+
\bar E_r^{\varepsilon,\mathrm{gr}}$. All norms in the next theorem
are taken before the collision limit.

\begin{theorem}[Complete optical source bounds at the collision scale]
\label{thm:v83-complete-weak-height}
For each previously fixed $r,L>0$ and $0<\eta<1/2$,
\begin{align}
 \limsup_m\sup_{R,n,k}\sup_{|J|=L}
       \mathcal N_{Q_\eta}(m^2D_r^\varepsilon;J)
                     &\le C_\eta(L+r),\label{eq:v83-weak-height}\\
 \limsup_m\sup_{R,n,k}\sup_{|J|=L}
       \mathcal N_{Q_\eta}(m^2D_{r,<b}^\varepsilon;J)
                     &\le C_\eta b_r^{\alpha-\eta}(L+r)
                      \quad(b_r<b_1).\label{eq:v83-critical-weak-height}
\end{align}
Here $\alpha=129/128$. For $1\le p<Q_\eta$, the left sides may be
replaced by the corresponding $L^p(J)$ norms, provided the right
sides are multiplied by
$(Q_\eta/(Q_\eta-p))^{1/p}L^{1/p-1/Q_\eta}$.
No essential-supremum conclusion is included.
\end{theorem}
\begin{proof}
Combine Theorem~\ref{thm:v83-potential-transport} with
Lemma~\ref{lem:v83-potential-norm}, then apply the fixed-width bounds
for the actual sources $v^\eta$ and $v^{\eta,b_r}$ in
Input~\ref{input:v83-local} on the interval of length $L+r$.
Its length and $b_r$ are fixed before the collision limit. The
auxiliary band is enlarged only afterwards.

If $\mathcal N_Q(f;J)\le A$, applying its definition to
$\{f>\lambda\}\cap J$ gives
$|\{f>\lambda\}\cap J|\le\min\{L,(A/\lambda)^Q\}$.
Integration of this bound against $p\lambda^{p-1}d\lambda$ proves
$\|f\|_{L^p(J)}\le (Q/(Q-p))^{1/p}A L^{1/p-1/Q}$.
Apply it at each finite count before taking the limits.
\end{proof}

For example, $\eta=1/4$ gives $Q_\eta=8/3$ and
$\alpha-\eta=97/128$. On intervals of length $r$, the complete
optical norm is $O(r)$; the part with initial $z<\sqrt r$ is
$O(r^{1+97/256})$. The estimates concern the untruncated original
source. They do not permit $p=\infty$, $\eta=1/2$, or a collision-
dependent roof width. Their role in the original pointwise problem
is made explicit in Section~\ref{sec:v83-pointwise}.
